\documentclass[oneside,a4paper,14pt]{extarticle}
\usepackage[a4paper,letterpaper,top=20mm,bottom=20mm,left=20mm,right=10mm]{geometry}
\usepackage[russian]{babel}
\usepackage{indentfirst}
\usepackage{graphicx}
\usepackage{caption}
\usepackage{titlesec}
\usepackage{minted, fancyvrb}
\usepackage{hyperref}
\usepackage{enumitem}
\usepackage{float}
\DeclareUnicodeCharacter{2500}{-}
\DeclareUnicodeCharacter{2550}{-}
\DeclareUnicodeCharacter{1F3CE}{} % 🏎
\DeclareUnicodeCharacter{1F3C6}{} % 🏆
\DeclareUnicodeCharacter{1F947}{} % 🥇
\DeclareUnicodeCharacter{26A1}{}  % ⚡
\DeclareUnicodeCharacter{1F4CB}{} % 📋
\DeclareUnicodeCharacter{1F3D7}{} % 🏗
\DeclareUnicodeCharacter{1F3C1}{} % 🏁
\DeclareUnicodeCharacter{1F468}{} % 🧑
\DeclareUnicodeCharacter{200D}{}  % ZWJ
\DeclareUnicodeCharacter{1F680}{} % 🚀
\DeclareUnicodeCharacter{270E}{}  % ✎


% Code styling
\setminted{style = rainbow_dash, fontsize = \small, breaklines}

% Header styles
\titleformat{\section}{\normalsize\bfseries}{\thesection}{1em}{}
\titleformat{\subsection}{\normalsize\bfseries}{\thesubsection}{1em}{}
\titleformat{\subsubsection}{\normalsize\bfseries}{\thesubsubsection}{1em}{}

% Spacing
\renewcommand\baselinestretch{1.33}
\setlength{\parindent}{1.25cm}

% List styles
\setlist[enumerate]{
  left=\parindent,
  label=\arabic*.,
  itemsep=0pt,
  topsep=5pt,
  partopsep=0pt,
  parsep=0pt
}

\setlist[itemize]{
  left=\parindent,
  itemsep=0pt,
  topsep=5pt,
  partopsep=0pt,
  parsep=0pt
}

% Hyperlinks
\hypersetup{
  colorlinks=true,
  linkcolor=black,
  urlcolor=blue,
  pdfborder={0 0 0},
  pdftitle={Система управления гоночными командами — Лабораторная работа 3},
  pdfauthor={Черкасов А.А.}
}

\begin{document}

\newpage
\thispagestyle{empty}
\begin{center}
  МИНИСТЕРСТВО НАУКИ И ВЫСШЕГО ОБРАЗОВАНИЯ РОССИЙСКОЙ ФЕДЕРАЦИИ ФЕДЕРАЛЬНОЕ ГОСУДАРСТВЕННОЕ БЮДЖЕТНОЕ ОБРАЗОВАТЕЛЬНОЕ УЧРЕЖДЕНИЕ ВЫСШЕГО ОБРАЗОВАНИЯ\\
  «ВЯТСКИЙ ГОСУДАРСТВЕННЫЙ УНИВЕРСИТЕТ»\\
  Институт математики и информационных систем\\
  Факультет автоматики и вычислительной техники\\
  Кафедра электронных вычислительных машин
\end{center}
\vspace{10mm}

\hfill
\begin{tabular}{l}
  \footnotesize Дата сдачи на проверку:                                          \\
  \footnotesize <<\rule[-1mm]{5mm}{0.10mm}\/>>\ \rule[-1mm]{20mm}{0.10mm}\ 2026 г. \\
  \footnotesize Проверено:                                                       \\
  \footnotesize <<\rule[-1mm]{5mm}{0.10mm}\/>>\ \rule[-1mm]{20mm}{0.10mm}\ 2026 г. \\
\end{tabular}
\vfill

\begin{center}
  Разработка приложения Formula Racing Management System на Java с использованием Spring Boot и PostgreSQL. \\
  Отчёт по лабораторной работе №3\\
  по дисциплине\\
  <<Объектно-ориентированное программирование>>\\
\end{center}
\vspace{25mm}
\noindent
\begin{tabular}{ll}
  Разработал студент гр. ИВТб-2301-05-00 & \hspace{18mm}\rule[-1mm]{30mm}{0.10mm}\,/Черкасов А. А./ \\
                                         & \hspace{25.5mm}\footnotesize(подпись)                    \\
  Преподаватель                  & \hspace{18mm}\rule[-1mm]{30mm}{0.10mm}\,/Шмакова Н. А./ \\
                                         & \hspace{25.5mm}\footnotesize(подпись)                    \\
\end{tabular}

\noindent
\begin{tabular}{lp{58mm}r}
  Работа защищена &  & \hspace{13mm}<<\rule[-1mm]{5mm}{0.10mm}\/>>\ \rule[-1mm]{30mm}{0.10mm}\ 2026 г.
\end{tabular}
\vfill

\begin{center}
  Киров\\
  2026
\end{center}

\newpage\thispagestyle{plain}

\section*{Цели лабораторной работы}
\begin{itemize}
  \item[$-$] освоить разработку веб-приложений на языке Java с использованием фреймворка Spring Boot;
  \item[$-$] изучить применение ORM-инструмента Hibernate для работы с реляционной базой данных PostgreSQL без написания SQL-запросов вручную;
  \item[$-$] реализовать полноценное CRUD"=приложение с веб-интерфейсом и нормализованной схемой БД;
  \item[$-$] изучить инструмент управления миграциями базы данных Flyway.
\end{itemize}

\section*{Задание}
Разработать веб-приложение <<Formula Racing Management System>> на языке Java, соответствующее следующим требованиям:
\begin{enumerate}
  \item Использовать язык программирования Java.
  \item Использовать СУБД PostgreSQL.
  \item Использовать сборщик Maven (зависимости только через \texttt{pom.xml}).
  \item Использовать Hibernate для работы с сущностями без SQL"=запросов.
  \item Нормализованная схема БД: не менее трёх связанных таблиц, структура соответствует 3НФ.
  \item Реализовать CRUD-операции: добавление, просмотр, редактирование и удаление объектов.
  \item Пользовательский интерфейс реализован в виде HTML"=страниц с использованием шаблонизатора Thymeleaf.
\end{enumerate}

\clearpage
\section*{Реализация приложения}

\subsection*{Предметная область}
Система автоматизирует управление гоночными командами в трёх чемпионатах: Formula 1, Formula 2 и Formula E. Каждая серия объединяет несколько команд, а каждая команда включает пилотов с личными очками. Иерархия объектов отражает реальные данные сезона 2026 года.

\subsection*{Архитектура приложения}
Проект построен по принципу многоуровневой архитектуры (Layered Architecture) в рамках Spring Boot:
\begin{itemize}
  \item[$-$] \textbf{Модели (Models)} --- сущности JPA, отражающие таблицы БД: \texttt{RacingSeries}, \texttt{RacingTeam}, \texttt{RacingDriver}.
  \item[$-$] \textbf{Репозитории (Repositories)} --- интерфейсы Spring Data JPA, автоматически предоставляющие CRUD-операции и производные запросы (например, \texttt{findByTeam\_SeriesIdOrderByPointsDesc}).
  \item[$-$] \textbf{Сервис (Service)} --- слой \texttt{RacingService}, инкапсулирующий бизнес-логику и изолирующий контроллер от деталей работы с данными.
  \item[$-$] \textbf{Контроллер (Controller)} --- \texttt{RacingController}, обрабатывающий HTTP-запросы и передающий данные в шаблоны Thymeleaf.
  \item[$-$] \textbf{Представление (View)} --- HTML"=шаблоны с разметкой Thymeleaf (\texttt{index.html}, \texttt{team-form.html}, \texttt{driver-form.html}).
\end{itemize}

\subsection*{Схема базы данных}
Для хранения данных используется PostgreSQL, развёрнутая в Docker-контейнере. Схема используется лабораторными работами №2, №3 и №4.

Схема нормализована до третьей нормальной формы (3НФ) и содержит три основные таблицы:

\begin{itemize}
  \item[$-$] \textbf{racing\_series}: хранит информацию о чемпионатах (название, официальное наименование, штаб-квартира, управляющая организация, год основания, количество этапов).
  \item[$-$] \textbf{racing\_teams}: хранит данные о командах (название, руководитель, база, силовая установка, очки). Таблица расширена дополнительными полями (победы, подиумы, бюджет и др.) для поддержки специфических полей GUI-приложения из ЛР №2. Каждая команда ссылается на серию через внешний ключ \texttt{series\_id}.
  \item[$-$] \textbf{racing\_drivers}: хранит данные о пилотах (имя, фамилия, гоночный номер, очки). Каждый пилот привязан к команде через \texttt{team\_id}.
\end{itemize}

Управление миграциями осуществляется через \textbf{Flyway}: при запуске приложения автоматически применяются миграции \texttt{V1}–\texttt{V5}. Скрипты создают таблицы, наполняют их реальными данными сезона 2026 года и расширяют структуру под потребности дочерних классов C\#-приложения.


\subsection*{JPA-сущности и связи}
Связи между сущностями реализованы через аннотации JPA:

\begin{minted}{java}
// В RacingSeries.java
@OneToMany(mappedBy = "series", cascade = CascadeType.ALL)
private List<RacingTeam> teams;

// В RacingTeam.java
@ManyToOne(fetch = FetchType.LAZY)
@JoinColumn(name = "series_id", nullable = false)
private RacingSeries series;

@OneToMany(mappedBy = "team", cascade = CascadeType.ALL)
private List<RacingDriver> drivers;
\end{minted}

Такой подход позволяет Hibernate автоматически выстраивать граф объектов при загрузке.

\subsection*{Оптимизация запросов: N+1 и @EntityGraph}
При отображении личного зачёта пилотов по сериям возникает проблема N+1 запроса: после загрузки списка пилотов Hibernate выполняет отдельный SQL-запрос для каждого связанного объекта \texttt{team} и \texttt{team.series}. Для решения используется \texttt{@EntityGraph} на производном запросе Spring Data:

\begin{minted}{java}
public interface RacingDriverRepository extends JpaRepository<RacingDriver, Long> {
    @EntityGraph(attributePaths = {"team", "team.series"})
    List<RacingDriver> findByTeam_SeriesIdOrderByPointsDesc(Long seriesId);
}
\end{minted}

Аннотация \texttt{@EntityGraph} указывает Hibernate загрузить связанные сущности одним LEFT JOIN запросом (используя \texttt{JOIN FETCH}), полностью устраняя N+1. Важно: \texttt{@EntityGraph} применяется только к конкретному методу, а не к \texttt{findAll()}, чтобы избежать ненужных JOIN-ов при простом перечислении всех пилотов.

Сервис формирует \texttt{LinkedHashMap<RacingSeries, List<RacingDriver>>}, который Thymeleaf итерирует для отрисовки таблицы личного зачёта с разделителями по сериям.

\subsection*{Веб-интерфейс}
Интерфейс реализован как многостраничное веб-приложение (7 шаблонов Thymeleaf).

Страницы приложения:
\begin{itemize}
  \item[$-$] \texttt{index.html} — дашборд с таблицей команд, пилотов и личным зачётом, сгруппированными по сериям;
  \item[$-$] \texttt{series.html} — просмотр/управление чемпионатами;
  \item[$-$] \texttt{teams.html} / \texttt{drivers.html} — списки команд и пилотов с фильтрацией;
  \item[$-$] \texttt{series-form.html} / \texttt{team-form.html} / \texttt{driver-form.html} — формы добавления/редактирования.
\end{itemize}

\section*{Вывод}
В ходе выполнения лабораторной работы было разработано веб-приложение на Java с использованием Spring Boot, PostgreSQL, Hibernate и Thymeleaf. Реализованы нормализованная схема БД в 3НФ, миграции Flyway, CRUD-операции для трёх сущностей, группировка данных по сериям и личный зачёт пилотов.

\newpage

\section*{Приложение А1. Сущность RacingSeries.java}
\inputminted{java}{code/src/main/java/com/vyatsu/racing/models/RacingSeries.java}

\section*{Приложение А2. Сущность RacingTeam.java}
\inputminted{java}{code/src/main/java/com/vyatsu/racing/models/RacingTeam.java}

\section*{Приложение А3. Сущность RacingDriver.java}
\inputminted{java}{code/src/main/java/com/vyatsu/racing/models/RacingDriver.java}

\section*{Приложение А4. RacingDriverRepository.java}
\inputminted{java}{code/src/main/java/com/vyatsu/racing/repositories/RacingDriverRepository.java}

\section*{Приложение А5. RacingService.java}
\inputminted{java}{code/src/main/java/com/vyatsu/racing/services/RacingService.java}

\section*{Приложение А6. RacingController.java}
\inputminted{java}{code/src/main/java/com/vyatsu/racing/controllers/RacingController.java}

\section*{Приложение А7. Скрипт миграции V1\_\_init.sql}
\inputminted{sql}{code/src/main/resources/db/migration/V1__init.sql}

\section*{Приложение А8. RacingTeamRepository.java}
\inputminted{java}{code/src/main/java/com/vyatsu/racing/repositories/RacingTeamRepository.java}

\section*{Приложение А9. RacingSeriesRepository.java}
\inputminted{java}{code/src/main/java/com/vyatsu/racing/repositories/RacingSeriesRepository.java}

\section*{Приложение Б1. Веб-интерфейс: Стили (static/style.css)}
\inputminted{css}{code/src/main/resources/static/style.css}

\section*{Приложение Б2. Главная страница (templates/index.html)}
\inputminted{html}{code/src/main/resources/templates/index.html}

\section*{Приложение Б3. Чемпионаты (templates/series.html)}
\inputminted{html}{code/src/main/resources/templates/series.html}

\section*{Приложение Б4. Команды (templates/teams.html)}
\inputminted{html}{code/src/main/resources/templates/teams.html}

\section*{Приложение Б5. Пилоты (templates/drivers.html)}
\inputminted{html}{code/src/main/resources/templates/drivers.html}

\section*{Приложение Б6. Форма чемпионата (templates/series-form.html)}
\inputminted{html}{code/src/main/resources/templates/series-form.html}

\section*{Приложение Б7. Форма команды (templates/team-form.html)}
\inputminted{html}{code/src/main/resources/templates/team-form.html}

\section*{Приложение Б8. Форма пилота (templates/driver-form.html)}
\inputminted{html}{code/src/main/resources/templates/driver-form.html}

\end{document}
